\documentclass[11pt]{article}
\usepackage[margin=1in,top=.85in,bottom=.85in]{geometry}
\usepackage{amsmath,amssymb,amsthm,hyperref}
\hypersetup{colorlinks=true,linkcolor=blue,urlcolor=blue,citecolor=blue}
\setlength{\emergencystretch}{2em}
\newtheorem{theorem}{Theorem}
\newtheorem{lemma}[theorem]{Lemma}
\title{A strict valley in Keith's multicolor partition array}
\author{Research proof and formalization packet}
\date{October 26, 2026}
\begin{document}
\maketitle
\begin{abstract}
Keith conjectured that every linear cut through a box-partition coefficient
array obtained by fixing one color variable to one is unimodal. We give a
three-color counterexample in the $25\times25$ box at partition size $325$.
Three consecutive points on a line of step $(1,-2)$ have coefficients
$53130$, $50600$, and at least $52900$, respectively. An injection from
three-element subsets of the allowed part sizes supplies the final bound.
A Lean certificate verifies the literal box-partition polynomial and its
failure of unimodality. The original two-color conjecture of Dousse and Kim
is not refuted.
\end{abstract}

\section{The source question}
A partition fits in an $M\times N$ box if it has at most $N$ parts and
every part is at most $M$. In Keith's $(k,1)$-coloring convention, each
distinct part size receives one of $k$ colors; a variable records how many
distinct sizes receive that color. Let $a_r$ count ordinary partitions of
a prescribed integer $n$ in the box having exactly $r$ distinct sizes.
For three colors, setting the third color variable to one gives
\begin{equation}\label{eq:F}
 F(x,y)=\sum_{r=0}^{\min(M,N)}a_r(1+x+y)^r.
\end{equation}
This is both the direct counting rule and the specialization of
Keith's Theorem~3 \cite{Keith}.

Conjecture~1 of \cite{Keith} says every linear cut through the remaining
coefficient array is unimodal. The discussion on printed page~8 explicitly
includes slant cuts with step $(1,-2)$. Thus the line used below is within
the stated intended scope. MathDB~345259 records this multicolor question.
It extends Dousse--Kim's two-color conjecture \cite[Conjecture~7.1(ii)]{DK};
we make no assertion against that original one-variable conjecture.

\begin{theorem}\label{thm:main}
Take $M=N=25$ and $n=325$, and define $C(i,j)=[x^iy^j]F(x,y)$ by
\eqref{eq:F}. Then
\[
 C(20,5)=53130,\qquad C(21,3)=50600,\qquad C(22,1)\ge52900.
\]
Consequently the full nonnegative-lattice cut $y=45-2x$ is not unimodal.
In particular, Keith's three-color linear-cut conjecture is false.
\end{theorem}

Here unimodal means weakly increasing up to a peak and weakly decreasing
thereafter. A strict valley at three consecutive points rules this out,
irrespective of the values elsewhere on the cut. All three coefficients
in Theorem~\ref{thm:main} are positive.

\section{Partitions near the maximum number of distinct sizes}
Throughout the proof the allowed sizes are $1,\ldots,25$ and their sum
is $325$. The staircase partition containing each size once is therefore
the unique partition with $25$ distinct sizes. Thus $a_{25}=1$.

There is no partition with $24$ distinct sizes. Indeed, let $b$ be the
omitted size. The distinct sizes sum to $325-b$, and the restriction to
at most $25$ parts permits at most one additional part. That part would
have to have size $b$, contradicting its omission. Hence $a_{24}=0$.

\begin{lemma}\label{lem:injection}
There are at least $\binom{25}{3}=2300$ partitions of $325$ in the box
with exactly $23$ distinct sizes.
\end{lemma}
\begin{proof}
Choose $1\le u<v<w\le25$, and put $d=u+w-v$. Then $u<d<w$.
Starting with the staircase, perform the following replacement:
\[
\begin{cases}
\text{remove $u,w$ and add $v,d$},&v\le d,\\
\text{remove $d,v$ and add $u,w$},&d<v.
\end{cases}
\]
The sum and number of parts are unchanged, and the added sizes are
disjoint from the removed sizes. The result has $25$ parts and exactly
$23$ distinct sizes. When $v=d$, the first case adds two further
copies of that size and still has exactly the claimed support.

The partition determines its two missing sizes and its multiset of two
extra parts, obtained by removing one occurrence of every present size.
If the smaller missing size is below the smaller extra size, the first
case applies: the missing sizes are $u,w$, and the smaller extra size is
$v$. Otherwise the second case applies: the extra sizes are $u,w$ and
the larger missing size is $v$. The original triple is therefore
uniquely recovered. This proves injectivity and the lemma.
\end{proof}

\section{The coefficient valley}
For a partition having $r$ distinct sizes, the number of color assignments
contributing to $x^iy^j$ after the third variable is set to one is
\begin{equation}\label{eq:weight}
 \binom ri\binom{r-i}{j}.
\end{equation}
Choose $i$ sizes for the first color, $j$ of the remaining sizes for the
second, and give the rest the third color. If $r<i+j$ the weight is zero.

At $(20,5)$ only $r=25$ contributes, giving $\binom{25}{5}=53130$.
At $(21,3)$ only $r=24,25$ could contribute; the former is absent, and
the staircase gives
\[
 \binom{25}{21}\binom43=50600.
\]
At $(22,1)$ each of the $2300$ distinct partitions from
Lemma~\ref{lem:injection} contributes
$\binom{23}{22}\binom11=23$. All other contributions are nonnegative, so
\[
 C(22,1)\ge23\binom{25}{3}=52900>50600.
\]
This proves Theorem~\ref{thm:main}. The three exponent pairs are
consecutive points on the indicated line. The full cut is indexed by
$x=0,\ldots,22$, so the nonunimodality is not produced by an arbitrary
reordering or omission of lattice points.

\section{Formal verification and attribution}
The certificate represents ordinary partitions with mathlib's
\texttt{Nat.Partition}: a multiset of positive natural numbers with
specified sum. It filters by the literal box restrictions, constructs
the polynomial \eqref{eq:F} as a sum over those partitions, and proves
its coefficient formula \eqref{eq:weight} using polynomial binomial
identities. No coefficient values or counting identities are endpoint
hypotheses.

For the injection, the kernel checks all $2300$ triples, their actual
partition multisets, box membership, support size and a decoding left
inverse. This finite certificate uses kernel reduction, not
\texttt{native\_decide}. The endpoint
\texttt{source\_full\_counterexample\_345259} states nonunimodality of
the literal polynomial cut. The six-module certificate has $37$
declarations; the only permitted axioms are \texttt{propext},
\texttt{Classical.choice} and \texttt{Quot.sound}. The separate proof
packet records pinned compiler/dependency revisions and replay receipts.

The source's counting formula, the two-color predecessor, binomial
coefficient identities and elementary finite-set counting are credited;
they are not separate discoveries. This is one proposed source-question
resolution. Initial searches found no earlier resolution of the
multicolor statement, but a finite search does not certify worldwide
priority. The actual arXiv source and its explicit slant-cut example were
inspected. The final publisher and institutional PDF requests returned
HTTP~403, so the final typeset body remains an explicit access limit.
Independent review and further bounded literature checks remain required
acceptance gates, recorded separately from this manuscript.

\begin{thebibliography}{9}
\bibitem{Keith} W.~J.~Keith, \emph{Restricted $k$-color partitions, II},
International Journal of Number Theory \textbf{17} (2021), no.~3,
591--601. \href{https://doi.org/10.1142/S1793042120400151}{doi:10.1142/S1793042120400151}.
\href{https://arxiv.org/abs/2001.08351v1}{arXiv:2001.08351v1},
Theorems~3--4, Conjecture~1 and printed page~8.
\bibitem{DK} J.~Dousse and B.~Kim, \emph{An overpartition analogue of
$q$-binomial coefficients, II: Combinatorial proofs and $(q,t)$-log
concavity}, Journal of Combinatorial Theory, Series~A \textbf{158}
(2018), 228--253.
\href{https://doi.org/10.1016/j.jcta.2018.03.011}{doi:10.1016/j.jcta.2018.03.011}.
\href{https://arxiv.org/abs/1701.07915}{arXiv:1701.07915}, Conjecture~7.1.
\end{thebibliography}
\end{document}
